\section{Discussion}\label{sec:discussion}

\subsection{Association, prediction and causation}

\paragraph{Association.} The three questions of the introduction have different answers.
The association is established: outside 2020--2021, faster GDP growth goes with a falling
unemployment rate in nearly every member state. Its size is not a single number:
\begin{itemize}
  \item it depends on whether the pandemic quarters are included;
  \item it depends on whether the error model absorbs residual persistence and a common
    European factor;
  \item the hierarchical sum is smaller than the classical one in \val{by_n_smaller} of
    \val{n_countries} countries, partly because the two estimate different quantities.
\end{itemize}

\paragraph{Prediction.} The association predicts little that the recent path of
unemployment does not already contain. The one decided comparison involves the quarter's
own GDP growth, which a forecaster would not know.

\paragraph{Causation.} No strategy available with public data identifies a causal Okun
coefficient for any shock. The coefficients reported here should be read as conditional
associations: summaries of how unemployment and output moved together in a period with a
particular mix of shocks.

\subsection{Positive, negative and unstable results}

\paragraph{Positive.} Results that the registered tests decide:
\begin{itemize}
  \item a negative association in \val{cl_n_neg_excluded} of \val{n_countries} countries
    without 2020--2021, robust to the core specifications in \val{cl_n_robust};
  \item \val{by_n_neg_all} sums credibly negative in every accepted Bayesian fit;
  \item Spain's large association, relative to the European mean;
  \item the Spain--Luxembourg and Spain--Netherlands differences, which survive every
    accepted fit and remain when labour input is measured by employment or hours;
  \item a more GDP-elastic output per hour in the Netherlands than in Spain, supported in
    the primary design and in the \val{cs_p22_settings} sensitivity settings;
  \item the conditional predictive gain over zero change before the pandemic
    (\val{pc_e1_zero}).
\end{itemize}

\paragraph{Negative.} Results where the registered test found nothing:
\begin{itemize}
  \item no decided gain from hierarchical pooling, GDP lags or the common factor in
    one-quarter-ahead forecasts;
  \item none of the \val{ss_published_holm} published post-2021 changes survives
    small-sample inference;
  \item no informative causal coefficient from external instruments, shift-share or sign
    restrictions;
  \item no sign that Luxembourg's commuters absorb its output fluctuations;
  \item \val{cs_supported} of \val{cs_total} case-study predictions supported.
\end{itemize}

\paragraph{Unstable.} Results that depend on choices or runs:
\begin{itemize}
  \item the magnitude of the association, through the pandemic quarters and the error
    model;
  \item \val{by_pairs_primary} credible country differences in the primary fit against
    \val{by_pairs_surviving} in every accepted fit;
  \item the Dutch delayed response, which holds in \val{claim_nl_holds} of
    \val{by_n_accepted} accepted fits;
  \item Portugal's significance across samples;
  \item two case-study predictions supported only before 2020;
  \item the employment-protection correlations, each decided in one of three variants;
  \item the acceptance of two Bayesian sensitivity fits, which differed between the
    Windows run and its Linux replication. The claim verdicts were identical under both
    acceptance sets, and the report states the robustness judgements under each.
\end{itemize}

\subsection{Limitations}

\paragraph{Data.}
\begin{itemize}
  \item All data are revised vintages, so the forecasts are retrospective.
  \item The Labour Force Survey's 2021 redesign introduces flagged and unflagged
    discontinuities. The Bayesian models mask the flagged observations, but the Dutch
    phase-in is not flagged.
  \item Monthly unemployment is a direct estimate in only eight countries.
  \item Hours per worker and employment in persons are unavailable for seven countries,
    including Portugal.
\end{itemize}

\paragraph{Models.}
\begin{itemize}
  \item The primary Bayesian model understates the persistence of unemployment changes.
    The registered amendment that removes this misfit has a common factor correlated with
    average GDP growth (\val{dyn_factor_corr}), so its country sums measure the association
    beyond part of the common GDP cycle.
  \item The post-2021 sample has about seventeen quarters. It is too short to test a
    change in the association with useful power.
  \item Country-level correlations with institutions use at most twenty countries and a
    single year of indicators.
\end{itemize}

\subsection{Open questions}

The study leaves these questions unresolved:
\begin{enumerate}
  \item Do countries differ because their labour markets respond differently to the same
    shock, or because they face different mixes of shocks?
    Equation~\eqref{eq:weighted} shows that both produce different Okun coefficients;
    separating them requires shock-specific identification.
  \item What is the effect on unemployment of a productivity or labour-supply shock in an
    EU country? This needs utilisation-adjusted productivity and quasi-experimental
    variation in labour supply.
  \item Why does output per hour absorb almost all output fluctuations in Luxembourg?
  \item Do temporary contracts carry Spain's employment adjustment? The direction is
    consistent with that view, but the registered test is not decided.
  \item Did the association weaken after 2021? This can only be answered with a longer
    sample.
\end{enumerate}
